% =====================================================================
% KHẢO LUẬN CỦA BẢN CHUYỂN NGỮ
% Không thuộc nguyên tác NCTM.
% Đề nghị đặt sau Chương 2, trước Phần II của FHSM.
% =====================================================================

\cleardoublepage
\gdef\CurrentRunningHead{Suy lí và kiến tạo ý nghĩa}

\bookmarksetup{startatroot}
\phantomsection
\addcontentsline{toc}{chapter}{Suy lí và kiến tạo ý nghĩa trong lớp học toán Việt Nam}

\begingroup
\clubpenalty=10000
\widowpenalty=10000
\displaywidowpenalty=10000

\thispagestyle{plain}

\begin{center}
{\small\bfseries KHẢO LUẬN CỦA BẢN CHUYỂN NGỮ\par}

\vspace{0.9\baselineskip}

{\Large\bfseries Suy lí và kiến tạo ý nghĩa\par}

\vspace{0.15\baselineskip}

{\Large\bfseries trong lớp học toán Việt Nam\par}
\end{center}

\vspace{1.15\baselineskip}

\begingroup
\small

\noindent
\textit{Ghi chú về tư cách văn bản.}
Khảo luận này do người chuyển ngữ biên soạn và không thuộc nguyên tác của Hội đồng Quốc gia Giáo viên Toán học Hoa Kì (National Council of Teachers of Mathematics, NCTM). Bài viết được đặt sau hai chương đầu của \emph{Focus in High School Mathematics: Reasoning and Sense Making} (FHSM), nơi NCTM đã trình bày các khái niệm suy lí (\emph{reasoning}) và kiến tạo ý nghĩa (\emph{sense making}) cùng những thói quen suy lí gắn với chúng.

\vspace{0.75\baselineskip}

\noindent
\textit{Tóm tắt.}
Khảo luận này diễn giải suy lí và kiến tạo ý nghĩa theo một cách gần với trải nghiệm học toán, rồi đặt hai quá trình ấy vào đối thoại với giáo dục toán học Việt Nam. Ví dụ về tổng ba góc của tam giác được dùng trước hết để làm nổi suy lí, từ quan sát và phỏng đoán đến việc xem xét giới hạn của chứng cứ thực nghiệm và tìm một căn cứ có tính tổng quát. Ví dụ về công thức khoảng cách trong mặt phẳng tọa độ tiếp đó làm nổi kiến tạo ý nghĩa, khi một công thức không còn đứng riêng lẻ mà được kết nối với độ chênh theo hai phương, tam giác vuông và định lí Pythagoras. Sau khi phân biệt hai quá trình theo cách ấy, bài viết chỉ ra rằng trong hoạt động toán học chúng thường đan xen và hỗ trợ lẫn nhau.

Trên nền đó, bài viết đặt FHSM trong đối thoại với Chương trình giáo dục phổ thông môn Toán 2018, giáo viên và học sinh Việt Nam. Chương trình đã xác định tư duy và lập luận toán học là một thành phần của năng lực toán học và mô tả nhiều biểu hiện cần đạt; FHSM không thay thế hay bổ sung một năng lực còn thiếu, mà cung cấp thêm một cách nhìn vào những cơ hội học tập qua đó học sinh thực sự được tham gia vào suy lí và kiến tạo ý nghĩa. Với giáo viên, trọng tâm chuyển sang việc những quyết định dạy học đang mở ra hay lấy mất các cơ hội ấy. Với học sinh, trọng tâm chuyển từ việc chỉ nhận dạng dạng bài, ghi nhớ công thức và tái hiện thủ tục sang việc kết nối, tìm căn cứ, kiểm tra kết luận, nhận ra cấu trúc và tái dựng kiến thức khi cần.

Qua ba phía gồm chương trình, giáo viên và học sinh, khảo luận tập trung vào một vấn đề chung. Trong quá trình học toán, học sinh có thực sự được tạo cơ hội để phát triển sự hiểu biết bằng cách kết nối tình huống, bối cảnh hoặc khái niệm đang xét với tri thức đã có, đồng thời suy lí từ chứng cứ hoặc giả định để đi tới những kết luận có căn cứ hay không? Câu hỏi ấy làm rõ giá trị mà FHSM có thể đem lại khi được đặt trong bối cảnh lớp học toán Việt Nam.

\endgroup

\clearpage

\section*{1. Một định lí quen thuộc và suy lí}

Học sinh được học rằng tổng ba góc của một tam giác bằng $180^\circ$.

Có em ghi nhớ ngay
$$
\widehat A+\widehat B+\widehat C=180^\circ.
$$
Sau đó, mỗi khi gặp bài toán cho hai góc, em lấy $180^\circ$ trừ đi hai góc ấy để tìm góc còn lại. Em làm đúng và có thể làm rất nhanh.

Nhưng một học sinh khác có thể hỏi \emph{tại sao lại là $180^\circ$?} Câu hỏi ấy mở ra một hoạt động toán học khác.

Giáo viên có thể cho học sinh vẽ nhiều tam giác, đo ba góc và cộng lại. Các kết quả đều gần $180^\circ$. Có thể cắt một tam giác bằng giấy, tách ba góc rồi ghép chúng cạnh nhau. Ba góc tạo thành một góc bẹt.

Từ những trường hợp quan sát được, học sinh có thể phỏng đoán rằng tổng ba góc của mọi tam giác đều bằng $180^\circ$. Nhưng ngay lúc ấy một câu hỏi khác xuất hiện. Làm như vậy với một vài tam giác thì làm sao biết kết luận đúng với mọi tam giác?

Đây là chỗ suy lí hiện ra rất rõ. Học sinh không chỉ có một kết quả hay một phỏng đoán, mà bắt đầu xem xét căn cứ của kết luận. Các em có thể thử một trăm tam giác, thậm chí một nghìn tam giác, nhưng những kết quả ấy vẫn chỉ là những trường hợp đã thử. Hơn nữa, phép đo luôn có sai số. Chứng cứ thực nghiệm có thể làm cho phỏng đoán trở nên đáng tin hơn, nhưng tự nó chưa cho phép khẳng định kết quả đúng với mọi tam giác.

Muốn đi xa hơn, ta cần một căn cứ không phụ thuộc vào việc tam giác được vẽ to hay nhỏ, nhọn hay tù. Qua một đỉnh của tam giác, ta vẽ một đường thẳng song song với cạnh đối diện. Những quan hệ góc tạo bởi hai đường thẳng song song cho phép chuyển hai góc còn lại của tam giác tới những vị trí kề với góc ở đỉnh ấy. Ba góc lúc này cùng tạo nên một góc bẹt. Vì thế tổng của chúng bằng $180^\circ$.

Trong câu chuyện này, suy lí trải từ việc hình thành phỏng đoán trên cơ sở quan sát đến việc nhận ra giới hạn của chứng cứ thực nghiệm, rồi tìm một lập luận cho phép rút ra kết luận có tính tổng quát. Chứng minh là phần chặt chẽ nhất của quá trình ấy, nhưng suy lí đã bắt đầu từ trước khi chứng minh xuất hiện.

Nếu chỉ nhìn ví dụ đến đây, người đọc hoàn toàn có thể cảm thấy rằng nói về suy lí đã đủ. Cảm giác đó có cơ sở, bởi ví dụ tổng ba góc của tam giác đặc biệt thích hợp để làm nổi câu hỏi FHSM đặt vào trung tâm của suy lí. Ta đang dựa vào đâu để đi tới kết luận, và căn cứ ấy cho phép kết luận đến mức nào?

Để thấy rõ hơn vì sao FHSM còn cần đến kiến tạo ý nghĩa, ta chuyển sang một ví dụ khác.

\section*{2. Kiến tạo ý nghĩa và sự đan xen với suy lí}

FHSM định nghĩa kiến tạo ý nghĩa là việc phát triển sự hiểu biết về một tình huống, bối cảnh hoặc khái niệm bằng cách kết nối nó với tri thức đã có. Nói gần với trải nghiệm học toán hơn, một đối tượng toán học trở nên có nghĩa hơn khi nó không còn đứng riêng lẻ mà được đặt vào những quan hệ mà người học đã hiểu và có thể huy động.

Công thức khoảng cách trong mặt phẳng tọa độ cho thấy điều này khá rõ. Giả sử học sinh cần tìm khoảng cách giữa hai điểm. Một cách học là nhớ
$$
d=\sqrt{(x_2-x_1)^2+(y_2-y_1)^2}.
$$
Nếu công thức chỉ được ghi nhớ như một chuỗi kí hiệu, học sinh có thể sử dụng đúng nhưng vẫn dễ lẫn nó với công thức khác hoặc quên mất một thành phần.

Một cách nhìn khác bắt đầu từ chính cấu trúc hình học của tình huống. Hiệu $x_2-x_1$ cho độ chênh theo phương ngang, còn $y_2-y_1$ cho độ chênh theo phương dọc. Hai độ chênh ấy tạo thành hai cạnh góc vuông của một tam giác vuông, trong khi khoảng cách cần tìm là cạnh huyền. Khi đó
$$
\Delta x=x_2-x_1,\quad \Delta y=y_2-y_1,
$$
và định lí Pythagoras cho
$$
d^2=(\Delta x)^2+(\Delta y)^2.
$$
Từ đó công thức khoảng cách có thể được dựng lại.

Điểm quan trọng ở đây không phải là ghi nhớ có giá trị thấp hơn việc tự dựng công thức. Một người làm toán hiệu quả phải ghi nhớ rất nhiều điều. Khác biệt nằm ở chỗ công thức có còn là một mảnh tri thức tương đối cô lập hay đã được kết nối với độ chênh tọa độ, tam giác vuông và định lí Pythagoras. Qua những kết nối ấy, người học phát triển sự hiểu biết về công thức và có một cấu trúc để quay về khi cần kiểm tra hoặc tái dựng nó. Đó là điều FHSM gọi là kiến tạo ý nghĩa.

Sự phân biệt với suy lí lúc này có thể thấy rõ hơn. Khi học sinh nhận ra công thức khoảng cách gắn với một tam giác vuông và với định lí Pythagoras, các em đang phát triển sự hiểu biết bằng cách kết nối điều đang xét với tri thức đã có. Khi từ những quan hệ ấy các em rút ra công thức khoảng cách, các em đang suy lí. Cùng một tình huống toán học có thể chứa cả hai quá trình, nhưng hai cách mô tả nhấn vào hai phương diện khác nhau của hoạt động.

Ta cũng có thể quay lại định lí tổng ba góc của tam giác. Việc đo nhiều tam giác, đánh giá sức mạnh của chứng cứ và xây dựng chứng minh làm suy lí hiện ra rõ nhất. Tuy nhiên, chứng minh bằng đường thẳng song song đồng thời có thể làm cho kết quả $180^\circ$ trở nên có nghĩa hơn. Ba góc vốn nằm ở ba đỉnh khác nhau được liên hệ với những góc kề nhau trên một đường thẳng; kết quả $180^\circ$ vì thế được nối với góc bẹt và những quan hệ góc mà học sinh đã biết. Chính một hoạt động chứng minh vừa có thể cung cấp căn cứ cho kết luận, vừa có thể làm sâu thêm sự hiểu biết về kết quả ấy.

Bởi vậy, suy lí và kiến tạo ý nghĩa không tạo thành hai bước cố định, cũng không chia lớp học thành hai loại hoạt động riêng. FHSM nhấn mạnh rằng chúng đan xen từ những quan sát phi hình thức đến những suy diễn hình thức. Có lúc những kết nối giúp người học biết nên suy lí theo hướng nào. Có lúc chính suy lí, một phản ví dụ hoặc một chứng minh làm thay đổi và làm sâu thêm cách hiểu.

Điều cần giữ lại không phải là cố gắn nhãn từng phút trong một bài học xem đâu là suy lí, đâu là kiến tạo ý nghĩa. Quan trọng hơn là nhận ra hai câu hỏi mà FHSM đặt ra. Người học đang phát triển sự hiểu biết bằng những kết nối nào? Và các em đang dựa vào chứng cứ hoặc giả định nào để đi tới những kết luận có căn cứ?

\section*{3. FHSM đối thoại với Chương trình môn Toán 2018}

Cuộc trao đổi không nên bắt đầu bằng việc hỏi Chương trình giáo dục phổ thông môn Toán 2018 có \emph{reasoning} hay chưa.

Chương trình đã xác định \emph{tư duy và lập luận toán học} là một trong năm thành phần của năng lực toán học, bên cạnh mô hình hóa toán học, giải quyết vấn đề toán học, giao tiếp toán học và sử dụng công cụ, phương tiện học toán. Chương trình cũng mô tả những biểu hiện của tư duy và lập luận như thực hiện các thao tác tư duy, sử dụng quy nạp và suy diễn, đưa ra chứng cứ và lí lẽ, giải thích, chứng minh, điều chỉnh cách giải quyết vấn đề.

Vì thế sẽ không đúng nếu nói Chương trình chỉ nêu mục tiêu chung chung, còn FHSM mới nói rõ học sinh phải làm gì. Chương trình đã nói khá nhiều về điều học sinh cần có khả năng thực hiện.

Một chuyên gia FHSM có lẽ sẽ đặt một câu hỏi khác. \emph{Những yêu cầu ấy sẽ trở thành trải nghiệm học toán thường xuyên của học sinh bằng cách nào?}

Đây là nơi hai văn bản có thể đối thoại.

Một chương trình quốc gia phải quy định mục tiêu, nội dung, yêu cầu cần đạt và cấu trúc năng lực. FHSM được viết cho một chức năng khác. Nó đi sâu vào việc suy lí và kiến tạo ý nghĩa phải hiện diện trong cách học những nội dung cụ thể như số, đại số, hàm số, hình học, thống kê và xác suất.

Vì vậy, không thật thích hợp nếu hỏi suy lí của FHSM rộng hơn hay hẹp hơn tư duy và lập luận toán học của Chương trình 2018. Hai thuật ngữ được đặt ở những vị trí khác nhau trong hai văn bản.

Trong Chương trình 2018, \emph{tư duy và lập luận toán học} là một thành phần của năng lực toán học cần được hình thành ở người học. Trong FHSM, \emph{suy lí} trước hết là một quá trình nền tảng mà người học phải thường xuyên tham gia khi làm toán. Nói cách khác, một bên đang mô tả một thành phần năng lực cần phát triển, còn bên kia nhấn mạnh một quá trình toán học mà học sinh cần thực hiện trong quá trình học. Hai khái niệm có nhiều điểm giao nhau, nhưng không nên vì thế mà đồng nhất chúng hoặc xếp chúng vào một quan hệ rộng hơn -- hẹp hơn.

Điều FHSM có thể đem vào cuộc đối thoại là một cách hỏi sâu hơn về \emph{cơ hội học tập}. Nếu chương trình yêu cầu học sinh biết lập luận và chứng minh, trong bao nhiêu bài học các em thực sự được tạo cơ hội hình thành phỏng đoán trước khi chứng minh? Nếu chương trình yêu cầu giải quyết vấn đề, học sinh có thường xuyên phải tự lựa chọn chiến lược, hay chiến lược đã được nhận diện sẵn thông qua ``dạng bài''? Nếu chương trình yêu cầu mô hình hóa, học sinh có phải cân nhắc giả thiết của mô hình và diễn giải kết quả trở lại tình huống ban đầu, hay chủ yếu thay số vào một mô hình đã được xây dựng sẵn? Nếu chương trình yêu cầu giao tiếp toán học, học sinh có phải nghe, hiểu và thẩm định cách nghĩ của người khác, hay giao tiếp chủ yếu được hiểu là trình bày lời giải đúng?

Những câu hỏi ấy không phủ nhận những gì chương trình đã làm. Chúng hướng sự chú ý từ một năng lực được xác định trong văn bản sang những trải nghiệm học tập cần có để năng lực ấy thực sự phát triển.

Kiến tạo ý nghĩa đặt thêm một câu hỏi đáng chú ý. Chương trình 2018 không có một thành phần năng lực mang tên ``kiến tạo ý nghĩa''. Tuy nhiên, những hoạt động có quan hệ chặt chẽ với nó xuất hiện ở nhiều nơi như kết nối các ý tưởng, chuyển đổi giữa những biểu diễn, mô hình hóa, giải quyết vấn đề, vận dụng và diễn giải.

FHSM không cho phép ta gom tất cả những điều ấy lại rồi tuyên bố rằng chúng ``thực ra là \emph{sense making}''. Nhưng khái niệm kiến tạo ý nghĩa giúp đặt một câu hỏi xuyên qua nhiều thành phần của chương trình. Khi học sinh học một nội dung, những tình huống, bối cảnh hoặc khái niệm đang xét được kết nối với tri thức đã có như thế nào?

Một chương trình có thể rất đầy đủ về nội dung nhưng học sinh vẫn có thể giữ các nội dung ấy thành những mảnh riêng biệt. Phương trình ở một ngăn, hàm số ở một ngăn, hình học ở một ngăn, xác suất ở một ngăn. Ngược lại, khi người học nhìn thấy những quan hệ xuyên qua các miền nội dung, tri thức bắt đầu trở thành một cấu trúc có thể huy động.

Theo nghĩa ấy, FHSM không phải một chương trình khác dành cho Việt Nam. Nó có thể đóng vai trò như một \emph{lăng kính bổ sung} giúp người thiết kế chương trình, sách giáo khoa và học liệu hỏi sâu hơn về những hoạt động toán học mà mỗi nội dung tạo cơ hội cho học sinh thực hiện.


\section*{4. FHSM đối thoại với giáo viên}

Với giáo viên, có lẽ điều đầu tiên một chuyên gia FHSM sẽ nói không phải là ``hãy dạy suy lí và kiến tạo ý nghĩa'', mà là đừng biến suy lí và kiến tạo ý nghĩa thành hai nội dung phải dạy thêm.

Nếu sau khi học xong đại số, hình học, thống kê lại cần thêm một tiết ``rèn suy lí'', thì ta đã hiểu sai tinh thần của FHSM. Suy lí và kiến tạo ý nghĩa phải xảy ra \emph{trong chính lúc học đại số, hình học, hàm số, thống kê và xác suất}.

Một câu hỏi thực tế hơn dành cho giáo viên là \emph{trong bài học này, phần việc trí tuệ quan trọng đang do ai thực hiện?}

Giả sử giáo viên muốn học sinh học công thức khoảng cách. Có thể trình bày công thức, giải một ví dụ rồi giao năm bài tương tự. Không có gì sai về mặt toán học. Học sinh vẫn cần luyện cách sử dụng công thức. Nhưng nếu đó là toàn bộ trải nghiệm, phần lớn cấu trúc khiến công thức có nghĩa đã được giáo viên làm thay.

Một cách tổ chức khác có thể bắt đầu từ hai điểm trên bản đồ hoặc hệ trục tọa độ. Giáo viên hỏi hai điểm cách nhau bao xa. Học sinh biết độ chênh ngang, độ chênh dọc, nhưng khoảng cách cần tìm lại là đoạn xiên. Từ đó xuất hiện tam giác vuông. Định lí Pythagoras được huy động. Công thức được hình thành. Sau đó giáo viên hoàn toàn có thể chuẩn hóa kí hiệu và cho học sinh luyện tập sử dụng.

Nội dung cuối cùng không đổi. Nhưng trải nghiệm dẫn tới nó khác. Ở trường hợp thứ hai, học sinh đã tham gia vào việc nhận ra cấu trúc làm cho công thức tồn tại.

Điều này cũng cho thấy FHSM không đối lập ``khám phá'' với ``giảng dạy trực tiếp''. Giáo viên không nhất thiết phải để học sinh tự phát minh mọi định nghĩa, định lí và phương pháp. Có những điều giáo viên nên nói. Có những kĩ thuật cần được chỉ dẫn. Có những thủ tục phải luyện tới mức thành thạo.

Câu hỏi quan trọng là sau những chỉ dẫn ấy, phần suy nghĩ nào vẫn còn dành cho học sinh. Giáo viên có thể yêu cầu học sinh nêu căn cứ cho một khẳng định, xét xem kết luận còn đúng khi thay giả thiết hay không, tìm phản ví dụ, so sánh ý tưởng chung của hai cách giải, giải thích vì sao một đáp số có vẻ vô lí, chuyển sang một biểu diễn khác hoặc tìm cách dựng lại một công thức đã quên.

Những câu hỏi như vậy không chỉ dùng để ``làm bài khó hơn''. Chúng thay đổi bản chất của việc học.

Hãy trở lại định lí tổng ba góc tam giác. Nếu giáo viên trình bày ngay chứng minh hoàn chỉnh và học sinh chép lại, lớp học có một chứng minh nhưng chưa chắc có nhiều suy lí từ phía học sinh. Nếu trước đó các em đã đo, phỏng đoán, tự hỏi việc đo có đủ hay không, rồi tìm kiếm một căn cứ có thể áp dụng cho mọi tam giác, chứng minh xuất hiện như lời giải cho một vấn đề mà chính học sinh đã nhận thấy.

Sự khác nhau không nằm trong chứng minh. Chứng minh có thể giống hệt nhau. Sự khác nhau nằm ở \emph{vai trò của chứng minh trong hoạt động nhận thức của người học}.

Một giáo viên vận dụng FHSM vì thế phải làm công việc khó hơn việc lựa chọn một ``phương pháp dạy học tích cực''. Người ấy phải liên tục phán đoán chỗ nào nên nói, chỗ nào nên hỏi, chỗ nào nên để học sinh thử, một sai lầm nào đáng khai thác, một lời giải chưa hoàn chỉnh nào đang chứa một ý tưởng tốt, khi nào cần yêu cầu thêm căn cứ, khi nào một giải thích hiện tại đã phù hợp với trình độ của học sinh và khi nào nên chuyển sang mức chặt chẽ cao hơn.

FHSM không cho giáo viên một thuật toán để quyết định những điều đó. Nó cung cấp một điểm nhìn để người dạy tự hỏi liệu quyết định của mình đang mở ra hay lấy mất cơ hội để học sinh trực tiếp làm công việc toán học.

\section*{5. FHSM đối thoại với học sinh}

Với học sinh, có lẽ không cần nói ngay về ``năng lực'', ``quá trình nhận thức'' hay ``kiến tạo ý nghĩa''. Có thể bắt đầu từ một thói quen rất quen. Gặp một bài toán mới, em thường hỏi điều gì trước tiên?

Nhiều học sinh sẽ hỏi \emph{bài này thuộc dạng nào?} Đó không phải một câu hỏi vô ích. Nhận ra cấu trúc quen thuộc là một phần quan trọng của năng lực toán học. Người làm toán giỏi cũng nhận ra những mẫu hình và cấu trúc mình từng gặp.

Vấn đề chỉ xuất hiện khi ``dạng bài'' trở thành điểm cuối của suy nghĩ. Học sinh thấy dạng, nhớ công thức, thay số rồi dừng lại.

Một người học toán ngày càng trưởng thành cần tập thêm những câu hỏi khác. Mình đã biết điều gì có liên quan tới bài này? Vì sao phương pháp này dùng được? Điều kiện nào của bài toán đang được sử dụng? Nếu thay một điều kiện thì chuyện gì xảy ra? Mình đang dự đoán điều gì? Một vài ví dụ có đủ để kết luận chưa? Có phản ví dụ không? Đáp số có hợp lí với tình huống không? Hai cách giải khác nhau có chung một cấu trúc nào? Nếu quên công thức, mình có thể tìm đường quay lại nó không?

Đây chính là chỗ suy lí và kiến tạo ý nghĩa bắt đầu trở thành \emph{thói quen của người học}, thay vì thuật ngữ trong một cuốn sách dành cho giáo viên.

Chẳng hạn, một học sinh nhớ công thức khoảng cách có thể giải bài rất nhanh. Điều đó tốt. Nhưng nếu một ngày quên công thức mà vẫn có thể nghĩ ``khoảng cách ngang và dọc tạo thành hai cạnh góc vuông; đoạn cần tìm là cạnh huyền'', rồi dựng lại công thức bằng định lí Pythagoras, em đang sở hữu kiến thức theo một cách mạnh hơn.

Tương tự, học sinh có thể biết chắc tổng ba góc tam giác là $180^\circ$. Nhưng nếu em còn phân biệt được việc đo một trăm tam giác chỉ tạo ra chứng cứ thực nghiệm với việc một chứng minh cho ta căn cứ để kết luận cho mọi tam giác trong hình học đang xét, thì em đã học thêm một điều rất quan trọng về chính bản chất của toán học.

Toán học không chỉ gồm những mệnh đề đúng. Nó còn có những chuẩn mực để quyết định \emph{vì sao ta được quyền nói một mệnh đề là đúng}.

Vì thế, một học sinh học theo tinh thần FHSM không nhất thiết phải chậm hơn, viết dài hơn hay lúc nào cũng giải thích mọi thao tác. Khi đã thành thạo, nhiều quá trình có thể trở nên nhanh và gần như tự động.

Điều quan trọng là phía sau sự thành thạo ấy vẫn còn một cấu trúc có thể được gọi trở lại khi cần. Học sinh biết mình đang làm gì, tại sao phương pháp có hiệu lực, điều kiện nào giới hạn nó và làm thế nào kiểm tra khi có điều bất thường.

Nói theo cách gần nhất với học sinh, học thuộc được là tốt, làm nhanh được cũng tốt. Nhưng hãy cố học đến mức nếu một ngày quên mất điều đã thuộc, em vẫn còn đường để tự tìm trở lại.

\section*{6. Một câu hỏi chung cho ba phía}

Nếu người thiết kế chương trình, giáo viên và học sinh cùng ngồi quanh một bàn, chuyên gia FHSM có lẽ sẽ không yêu cầu cả ba học thêm một bộ thuật ngữ mới. Người ấy có thể đặt cùng một vấn đề ở ba cấp độ.

Với \emph{người thiết kế chương trình}, câu hỏi là những mục tiêu năng lực sẽ được chuyển thành những cơ hội học tập nào để học sinh thực sự được tham gia vào suy lí và kiến tạo ý nghĩa.

Với \emph{giáo viên}, câu hỏi là những quyết định dạy học đang mở ra hay lấy mất những cơ hội để học sinh kết nối, hình thành sự hiểu biết, tìm căn cứ, rút ra và kiểm tra kết luận.

Với \emph{học sinh}, câu hỏi là khi làm toán, em chỉ đang tìm cách có đáp án, hay còn nhận ra điều đang học liên hệ với những gì đã biết, mình đang dựa vào đâu và vì sao kết luận có thể được chấp nhận.

Ba câu hỏi ấy khác nhau nhưng đều quy về việc học sinh có thực sự được tham gia vào suy lí và kiến tạo ý nghĩa trong quá trình học toán hay không.

Đó có lẽ là đóng góp quan trọng nhất mà FHSM có thể đem vào cuộc đối thoại với giáo dục toán học Việt Nam. Chương trình có thể xác định rất đúng năng lực cần đạt, sách giáo khoa có thể lựa chọn nội dung tốt, giáo viên có thể giảng chính xác và học sinh có thể hoàn thành bài tập với đáp số đúng. Nhưng những điều ấy chưa tự động cho biết học sinh đã được kết nối điều đang học với tri thức đã có như thế nào, đã dựa vào những căn cứ nào để hình thành và kiểm tra kết luận, hay hai quá trình ấy đã thực sự trở thành phần việc của chính người học đến đâu.

FHSM hướng sự chú ý vào khoảng không thường khó thấy ấy.

\section*{7. Suy lí và kiến tạo ý nghĩa như một cách nhìn việc học toán}

Sau khi đọc FHSM, có thể vẫn hỏi rốt cuộc suy lí và kiến tạo ý nghĩa khác nhau ở đâu. Có thể trả lời ngắn gọn như sau. Kiến tạo ý nghĩa là quá trình phát triển sự hiểu biết về một tình huống, bối cảnh hoặc khái niệm bằng cách kết nối nó với tri thức đã có; suy lí là quá trình rút ra kết luận trên cơ sở chứng cứ hoặc những giả định đã nêu. Vì thế, kiến tạo ý nghĩa đặt trọng tâm vào việc những gì đang được học trở nên có nghĩa trong mạng lưới hiểu biết của người học như thế nào, còn suy lí đặt trọng tâm vào căn cứ cho phép người học đi tới một kết luận. Tuy nhiên, trong hoạt động toán học thực tế, hai quá trình này thường đan xen. Những kết nối được hình thành trong quá trình kiến tạo ý nghĩa có thể cung cấp cơ sở cho suy lí, còn suy lí có thể làm thay đổi hoặc làm sâu thêm sự hiểu biết về tình huống, bối cảnh hay khái niệm đang xét.

Sự phân biệt ấy là cần thiết, nhưng không phải là điểm cuối của FHSM. Nếu chỉ cố vạch một đường biên thật sắc để xác định từng hành động của học sinh thuộc về suy lí hay kiến tạo ý nghĩa, ta có thể bỏ lỡ điều lớn hơn mà NCTM muốn làm.

Suy lí và kiến tạo ý nghĩa không được đưa ra để giáo viên có thêm hai ô trong kế hoạch bài dạy. Chúng cung cấp một cách nhìn vào việc học toán.

Ta vẫn cần hỏi nội dung nào phải được học, học sinh có sử dụng đúng khái niệm và thủ tục không, các em có đạt yêu cầu của chương trình không. Nhưng FHSM đòi hỏi thêm một lớp câu hỏi. Điều đang học được kết nối với tri thức đã có ra sao? Những kết nối ấy đang làm sự hiểu biết của học sinh phát triển như thế nào? Các em hình thành kết luận từ những căn cứ nào? Đã kiểm tra phỏng đoán thế nào? Khi gặp một kết quả trái với dự đoán, các em có sửa cách nghĩ không? Một chứng minh chỉ được chép lại hay còn giúp làm sáng cấu trúc của vấn đề? Một thủ tục chỉ được ghi nhớ hay còn có thể được giải thích, kiểm tra và tái dựng?

Nếu những câu hỏi ấy được đặt thường xuyên, suy lí và kiến tạo ý nghĩa không còn là hai thuật ngữ đứng bên ngoài chương trình. Chúng hiện diện trong cách học sinh học một phương trình, dựng một công thức, đọc một đồ thị, xây dựng một chứng minh, lựa chọn một mô hình, diễn giải một bảng số liệu hay kiểm tra một kết quả.

Đó cũng là lí do FHSM không cần được tiếp nhận ở Việt Nam như một mô hình để sao chép. Giá trị của nó nằm ở một câu hỏi giản dị nhưng khó trả lời.

\begin{quote}
\emph{Trong khi học một nội dung toán học, học sinh có được tham gia vào những hoạt động mà ở đó việc phát triển sự hiểu biết bằng cách kết nối tình huống, bối cảnh hoặc khái niệm đang xét với tri thức đã có đan xen với việc rút ra kết luận trên cơ sở chứng cứ hoặc giả định đã nêu hay không?}
\end{quote}

Một chương trình hướng tới năng lực cần quan tâm đến những cơ hội học tập ấy. Giáo viên phải tổ chức để chúng thực sự xuất hiện trong bài học. Và học sinh chỉ có thể phát triển suy lí và kiến tạo ý nghĩa khi chính các em được tham gia vào những hoạt động đó.

\endgroup
